\SourcePage{142}
\SGAchapter{XVI}{光滑态射的基变换定理及其应用}{Théorème de changement de base par un morphisme lisse, et applications}
\begin{center}M. Artin\end{center}
\SGAsection{1}{光滑态射的基变换定理}
\begin{SGAresult}{定理}{1.1}
设 $L\subset\mathbb P$，$n\in\mathbb N$，且
\[
\begin{tikzcd}
X\arrow[d,"f" ]&X'\arrow[l,"g'"']\arrow[d,"f'"]\\
Y&Y'\arrow[l,"g"]
\end{tikzcd}\tag{*}
\]
为\Term{cartesian-square}{笛卡尔图}，其中 $f$ \Term{quasi-compact}{拟紧}且\Term{quasi-separated}{拟分离}，$g$ 是\Term{universally-locally-acyclic}{普遍局部 $0$-无环}的（相应地，对 $L$ \Term{universally-locally-aspherical}{普遍局部 $1$-无球}；相应地，对 $L$ \Term{universally-locally-acyclic}{普遍局部 $n$-无环}）（第 XV 讲 1.11）。则对 $X$ 上每个\Term{set}{集合}\Term{sheaf}{层} $F$（相应地，$\operatorname{ind}$-$L$-群\Term{sheaf}{层}；相应地，$L$-挠阿贝尔\Term{sheaf}{层}），\Term{base-change}{基变换}\Term{morphism}{态射}（第 XII 讲 4.1.2）
\[
\varphi^q:g^*(R^qf_*)F\longrightarrow(R^qf'_*)g'^*F
\]
在 $q=0$ 时（相应地，$q\leq1$ 时；相应地，$q\leq n$ 时）为\Term{bijection}{双射}。
\end{SGAresult}
\SourcePage{143}
特别地，应用第 XV 讲 2.1，得到：
\begin{SGAresult}{推论}{1.2}
（\Term{smooth-morphism}{光滑态射}的\Term{base-change-theorem}{基变换定理}。）设 $(*)$ 中的\Term{morphism}{态射} $g$ \Term{smooth}{光滑}，$f$ \Term{quasi-compact}{拟紧}且\Term{quasi-separated}{拟分离}，$L\subset\mathbb P$ 为 $Y$ 的\Term{residue-characteristic}{剩余特征}\Term{set}{集合}的补集。则对 $X$ 上每个\Term{set}{集合}\Term{sheaf}{层} $F$（相应地，$\operatorname{ind}$-$L$-群\Term{sheaf}{层}；相应地，$L$-挠阿贝尔\Term{sheaf}{层}），上述\Term{base-change}{基变换}\Term{morphism}{态射} $\varphi^q$ 在 $q=0$ 时（相应地，$q=0,1$ 时；相应地，对每个 $q$）为\Term{bijection}{双射}。
\end{SGAresult}
\begin{proof}[1.1 的证明]
先处理 $f:X\to Y$ \Term{quasi-projective}{拟射影}的情形：这只是把\Term{proper}{固有}\Term{morphism}{态射}的\Term{base-change}{基变换}定理（第 XII 讲 5.1）与第 XV 讲 1.9 的定义结合起来。事实上，有\Term{commutative-diagram}{交换图}
\[
\begin{tikzcd}
X\arrow[r,"i"]\arrow[dr,"f"']&\overline X\arrow[d,"\overline f"]\\
&Y
\end{tikzcd}
\]
其中 $\overline f$ \Term{projective}{射影}，因而\Term{proper}{固有}，$i$ 为\Term{open-immersion}{开浸入}。由第 XII 讲 5.1 知，无论 $g:Y'\to Y$ 如何，$\overline f$ 都满足\Term{base-change}{基变换}定理。由于 $g\times_Y\overline X:Y'\times_Y\overline X\to\overline X$ 仍是\Term{universally-locally-acyclic}{普遍局部 $0$-无环}的（相应地，……），由第 XII 讲 4.4 (ii)，立即归结为证明\Term{morphism}{态射} $i$ 的定理，也就是归结到 $f$ 为\Term{open-immersion}{开浸入}的情形。这时定理是第 XV 讲 1.10 (ii) 的定义的推论。

现在处理一般情形。断言在 $Y$ 上是\Term{local-scheme}{局部}的，故可设 $Y$ \Term{affine}{仿射}。
\end{proof}
\begin{SGAresult}{引理}{1.3}
还可设 $X$ \Term{affine}{仿射}。
\end{SGAresult}
\begin{proof}
如第 XII 讲 6.1 那样，为对给定的 $f$ 证明 1.1，只须对由 $Y',Y$ 的\Term{strict-localization}{严格局部化}之间、由 $g$ 诱导的\Term{local-morphism}{局部态射} $\widetilde g:\widetilde Y'\to\widetilde Y$，经\Term{base-change}{基变换} $\widetilde Y\to Y$、$\widetilde Y'\to Y'$ 得到的情形证明它；亦即证明\Term{homomorphism}{同态}
\[
H^q(\widetilde X,F)\longrightarrow H^q(\widetilde X',F'),\qquad F'=g'^*(F),\tag{*}
\]
对 $q\leq n$ 为\Term{bijection}{双射}，其中 $\widetilde X=X\times_Y\widetilde Y$，$\widetilde X'=X'\times_{Y'}\widetilde Y'$。因此，归结到 $Y,Y',g$ 均为\Term{strictly-local}{严格局部}的情形，并在承认对\Term{affine}{仿射} $f$ 已证明 1.1 的前提下，证明上述映射为\Term{bijection}{双射}。取 $X$ 的\Term{finite}{有限}\Term{affine}{仿射}\Term{open-cover}{开覆盖} $(X_i)$，令 $Z$ 为各 $X_i$ 的和\Term{scheme}{概形}，则 $Z$ \Term{affine}{仿射}，并带有\Term{surjective-morphism}{满射态射} $h:Z\to X$，从而有 $h':Z'\to X'$。注意 $h$ \Term{separated}{分离}，且当 $f$ \Term{separated}{分离}时 $h$ \Term{affine}{仿射}。于是应用\Term{descent}{下降}引理（第 XII 讲 6.8），为对 $q\leq n$ 及一切 $F$ 证明 $(*)$ 为\Term{bijection}{双射}，只须证明相应映射
\SourcePage{144}
\[
H^q(Z,h^*F)\longrightarrow H^q(Z',h'^*F')
\]
为\Term{bijection}{双射}，条件是 $Z\to X$ 关于\Term{base-change}{基变换} $X'\to X$ 的\Term{homomorphism}{同态}在\Term{dimension}{维数} $\leq n$ 为\Term{isomorphism}{同构}。如果 $f$ \Term{separated}{分离}，则 $h$ \Term{affine}{仿射}，由假设这确实成立，故当 $f$ \Term{separated}{分离}时 1.1 成立。一般情形下，再对始终\Term{separated}{分离}的 $h$ 应用刚才的结果，仍得 1.1 对 $f$ 成立，故得 1.3。
\end{proof}
\textbf{注记。}也可使用 $X$ 的\Term{open-cover}{开覆盖} $(X_i)$ 的 Leray \Term{spectral-sequence}{谱序列}及其非交换变体；但这种方法与调用第 XII 讲 6.8 本质上无异，却有一个缺点：它迫使我们重新区分惯常的三种情形，而该处已经做过这些工作。

\begin{proof}[1.1 的证明（续）]
为完成一般情形的证明，只须归结到 $f:X\to Y$ \Term{affine}{仿射}且\Term{finite-type}{有限型}、因而\Term{quasi-projective}{拟射影}的情形。由 1.3，可设 $X$ \Term{affine}{仿射}。令 $X=\varprojlim X_\alpha$，其中 $f_\alpha:X_\alpha\to Y$ 是 $Y$ 上\Term{affine}{仿射}且\Term{finite-type}{有限型}的\Term{scheme}{概形}。用通常的取\Term{limit}{极限}技术（第 VII 讲 5）即可完成证明。
\end{proof}
\textbf{补充 1.4。}a) 请读者自行观察，同样的证明在假设 $g$ \Term{universally-locally-acyclic}{普遍局部 $(-1)$-无环}（第 XV 讲 1.14）时仍给出一个结论：此时，对 $X$ 上任意\Term{set}{集合}\Term{sheaf}{层} $F$，\Term{base-change}{基变换}\Term{homomorphism}{同态}
\[
\phi^0:g^*(f_*F)\longrightarrow f'_*(g'^*F)
\]
是\Term{injectivity}{单射}。另一方面，重新设 $n\geq0$，还可用以下\Term{injectivity}{单射}性陈述补充 1.1，它也包含于上述证明中：若 $g$ \Term{universally-locally-acyclic}{普遍局部 $0$-无环}，则对 $X$ 上一切群\Term{sheaf}{层} $F$，1.1 的\Term{base-change}{基变换}\Term{homomorphism}{同态} $\phi^1$ 为\Term{monomorphism}{单态射}；若 $g$ 对 $L$ \Term{universally-locally-acyclic}{普遍局部 $n$-无环}，则对 $X$ 上任意 $L$-挠阿贝尔\Term{sheaf}{层} $F$，1.1 的\Term{homomorphism}{同态} $\phi^{n+1}$ 为\Term{monomorphism}{单态射}。

b) 从 $g$ 对 $L$ 的\Term{locally-aspherical}{局部 $1$-无球性}假设出发，引入 Giraud 学位论文的非交换 $2$-\Term{cohomological-invariant}{上同调不变量}，大概也能给出类似的\Term{injectivity}{单射}性结论，参见第 XII 讲 5.11\footnote{这正是 J. Giraud 的书中实际建立的结果。}；下面的 (1.6) 和 (2.3) 也有同样的问题。
\begin{SGAresult}{推论}{1.5}
设 $K/k$ 为域\Term{group-extension}{扩张}，$L=\mathbb P-\{p\}$，其中 $p=\operatorname{char}k$。则 $\operatorname{Spec}K\to\operatorname{Spec}k$ 对 $L$ \Term{universally-locally-acyclic}{普遍局部无环}，且对 $L$ \Term{universally-locally-aspherical}{普遍局部 $1$-无球}；如果 $k,K$ 均\Term{separably-closed}{可分闭}，则上述\Term{morphism}{态射}还对 $L$ \Term{universally-acyclic}{普遍无环}，且对 $L$ \Term{universally-aspherical}{普遍 $1$-无球}。
\end{SGAresult}
\SourcePage{145}
立即指出 (1.5) 的如下推论：
\begin{SGAresult}{推论}{1.6}
设 $K/k$ 为\Term{separably-closed}{可分闭}域的\Term{group-extension}{扩张}，$X$ 为\Term{scheme}{概形}，$F$ 为 $X$ 上的\Term{set}{集合}\Term{sheaf}{层}（相应地，$\operatorname{ind}$-$L$-群\Term{sheaf}{层}；相应地，$L$-挠\Term{abelian-group}{阿贝尔群}\Term{sheaf}{层}），其中 $L=\mathbb P-\{p\}$，$p=\operatorname{char}k$。则映射
\[
H^q(X,F)\longrightarrow H^q(X_K,F_K)
\]
在 $q=0$ 时（相应地，$q\leq1$ 时；相应地，对一切 $q$）为\Term{bijection}{双射}。
\end{SGAresult}
事实上，这立即由 $\operatorname{Spec}K\to\operatorname{Spec}k$ 对 $L$ \Term{universally-acyclic}{普遍无环}且\Term{universally-aspherical}{普遍 $1$-无球}及定义而得。注意，利用第 VIII 讲 2.4，也可把上述 $H^q$ 间的\Term{homomorphism}{同态}解释为\Term{base-change}{基变换}\Term{homomorphism}{同态}
\[
g^*(R^qf_*F)\longrightarrow R^qf'_*(g'^*F),
\]
其中 $f:X\to\operatorname{Spec}k$、$f':X_K\to\operatorname{Spec}K$、$g:\operatorname{Spec}K\to\operatorname{Spec}k$ 及 $g':X_K\to X$ 均为显然的\Term{morphism}{态射}。因此，当 $X$ \Term{quasi-compact}{拟紧}且\Term{quasi-separated}{拟分离}时，由 1.1，这些\Term{homomorphism}{同态}的\Term{bijection}{双射}性也可看作 $\operatorname{Spec}K\to\operatorname{Spec}k$ 对 $L$ \Term{universally-locally-acyclic}{普遍局部无环}且\Term{universally-locally-aspherical}{普遍局部 $1$-无球}的推论。后一事实已经蕴含\Term{quasi-compact}{拟紧}且\Term{quasi-separated}{拟分离} $X$ 的 1.6，而这显然足以推出 $\operatorname{Spec}K\to\operatorname{Spec}k$ 对 $L$ \Term{universally-acyclic}{普遍无环}且\Term{universally-aspherical}{普遍 $1$-无球}（第 XV 讲 1.7、1.6 (i)）。

这说明，为证明 1.5，只须证明它的第一个断言。由第 VIII 讲 1.1，可将 $k$ 换成其\Term{perfect-closure}{完美闭包}，故可设 $k$ \Term{perfect-field}{完美}。$k$ 的任意\Term{group-extension}{扩张} $K$ 都是其\Term{finite-type}{有限型} $k$-子\Term{algebra}{代数} $A$ 的\Term{inductive-limit}{归纳极限}，而 $k$ \Term{perfect}{完美}，故将这样的 $A$ 适当\Term{localization}{局部化}后，可以设它\Term{smooth}{光滑}，从而 $K$ 是\Term{smooth}{光滑}子\Term{algebra}{代数}的\Term{filtered-colimit}{滤过归纳极限}。根据第 XV 讲 2.1，这些子\Term{algebra}{代数}对 $L$ \Term{universally-locally-acyclic}{普遍局部无环}且\Term{universally-locally-aspherical}{普遍局部 $1$-无球}，再由第 XV 讲 1.13 (ii) 即得结论。\hfill$\square$

\textbf{注记 1.7。}可将不变性定理 1.5 与第 XII 讲 5.4 比较：后者不必假设所考虑的挠\Term{sheaf}{层}与\Term{residue-characteristic}{剩余特征}互素，但必须假设 $X$ 在 $k$ 上\Term{proper}{固有}。即使对 $H^1(X,\mathbb Z/p\mathbb Z)$，其中 $X$ 是\Term{characteristic}{特征} $p>0$ 的\Term{algebraically-closed}{代数闭}域 $k$ 上的\Term{affine}{仿射}\Term{algebraic-curve}{代数曲线}（例如\Term{affine}{仿射}直线），上述陈述的类似\Term{form}{形式}也不成立，原因在于无穷远处的“\Term{wild-ramification}{野分歧}”现象；它使这种情形下群 $\mathbb Z/p\mathbb Z$ 的主\Term{etale-covering}{平展覆盖}的分类本质上是“\Term{continuous}{连续}的”。
\SourcePage{146}
\SGAsection{2}{上同调群的特殊化定理}
\begin{SGAresult}{定理}{2.1}
设 $L\subset\mathbb P$，$n\in\mathbb N$，$f:X\to Y$ 为\Term{finite-presentation}{有限表示}的\Term{proper}{固有}\footnote{关于较\Term{properness}{固有性}弱而仍可得到相同结论的假设，参见 SGA 5 第 II 讲 3 和 SGA 1 第 XIII 讲。}\Term{morphism}{态射}，且为\Term{locally-acyclic}{局部 $0$-无环}的（相应地，对 $L$ \Term{locally-aspherical}{局部 $1$-无球}；相应地，对 $L$ \Term{locally-acyclic}{局部 $n$-无环}）。设 $F$ 是 $X$ 上\Term{constructible}{可构造}且\Term{locally-constant}{局部常值}的\Term{set}{集合}\Term{sheaf}{层}（相应地，$\operatorname{ind}$-$L$-群\Term{sheaf}{层}；相应地，$L$-挠阿贝尔\Term{sheaf}{层}）。则 $R^qf_*F$ 在 $q=0$ 时（相应地，$q=0,1$ 时；相应地，$q\leq n$ 时）\Term{constructible}{可构造}且\Term{locally-constant}{局部常值}，并且对 $Y$ 的任意\Term{geometric-point}{几何点} $\bar y$，对同样的 $q$ 值有
\[
(R^qf_*F)_{\bar y}\simeq H^q(X_{\bar y},F_{\bar y}).
\]
\end{SGAresult}
\textbf{注记。}最后一个断言无非是\Term{proper}{固有}\Term{morphism}{态射}的\Term{base-change}{基变换}定理（第 XII 讲 5.2）。
由第 XV 讲 2.1，立即得到如下推论：
\begin{SGAresult}{推论}{2.2}
（\Term{cohomology-group}{上同调群}的\Term{specialization}{特殊化}定理。）设 $f:X\to Y$ 为\Term{proper}{固有}\Term{smooth}{光滑}\Term{morphism}{态射}，$L\subset\mathbb P$ 为 $Y$ 的\Term{residue-characteristic}{剩余特征}\Term{set}{集合}的补集。设 $F$ 为 $X$ 上\Term{constructible}{可构造}的\Term{set}{集合}\Term{sheaf}{层}（相应地，$\operatorname{ind}$-$L$-群\Term{sheaf}{层}；相应地，$L$-挠阿贝尔\Term{sheaf}{层}）。则 $R^qf_*F$ 在 $q=0$ 时（相应地，$q=0,1$ 时；相应地，对一切 $q$）\Term{constructible}{可构造}且\Term{locally-constant}{局部常值}，并且对 $Y$ 的任意\Term{geometric-point}{几何点} $\bar y$，对同样的 $q$ 值有
\[
(R^qf_*F)_{\bar y}\simeq H^q(X_{\bar y},F_{\bar y}).
\]
\end{SGAresult}
\textbf{注记。}取 $n=1$ 并译成\Term{fundamental-group}{基本群}的语言，就重新得到比较\Term{proper}{固有}\Term{smooth}{光滑}\Term{morphism}{态射}的各\Term{fiber}{纤维}的\Term{fundamental-group}{基本群}的 SGA 1 第 X 讲 3.8。
\begin{proof}[2.1 的证明]
已知 $R^qf_*F$ \Term{constructible}{可构造}（第 XIV 讲 1.1），故由第 IX 讲 2.11，只须证明：对 $Y$ 的\Term{geometric-point}{几何点}间任意\Term{specialization}{特殊化}\Term{morphism}{态射} $\bar y_1\to\bar y_0$，相应的\Term{specialization}{特殊化}\Term{morphism}{态射}（第 VIII 讲 7.7）
\[
(R^qf_*F)_{\bar y_0}\longrightarrow(R^qf_*F)_{\bar y_1}\tag{*}
\]
对所考虑的 $q$ 值为\Term{bijection}{双射}。关于这一点，有如下稍广的结果（无须假设 $f$ \Term{finite-presentation}{有限表示}或 $F$ \Term{constructible}{可构造}）：
\end{proof}
\begin{SGAresult}{推论}{2.3}
设 $L\subset\mathbb P$，$n\in\mathbb N$，$f:X\to Y$ 为\Term{proper}{固有}且\Term{locally-acyclic}{局部 $0$-无环}的\Term{morphism}{态射}（相应地，对 $L$ \Term{locally-aspherical}{局部 $1$-无球}；相应地，对 $L$ \Term{locally-acyclic}{局部 $1$-无环}），$F$ 为 $X$ 上\Term{locally-constant}{局部常值}的\Term{set}{集合}\Term{sheaf}{层}（相应地，$\operatorname{ind}$-$L$-群\Term{sheaf}{层}；相应地，$L$-挠阿贝尔\Term{sheaf}{层}）。则对 $Y$ 的\Term{geometric-point}{几何点}间任意\Term{specialization}{特殊化}\Term{morphism}{态射} $\bar y_1\to\bar y_0$，
\SourcePage{147}
上述相应的\Term{specialization}{特殊化}\Term{homomorphism}{同态} $(*)$ 在 $q=0$ 时（相应地，$q=0,1$ 时；相应地，$q\leq n$ 时）为\Term{bijection}{双射}。此外，当 $n=0$，即 $g$ \Term{locally-acyclic}{局部 $0$-无环}时，对 $X$ 上任意\Term{locally-constant}{局部常值}群\Term{sheaf}{层} $F$，\Term{homomorphism}{同态}
\[
(R^1f_*F)_{\bar y_0}\longrightarrow(R^1f_*F)_{\bar y_1}
\]
是\Term{injectivity}{单射}；对任意 $n$ 及 $X$ 上任意\Term{locally-constant}{局部常值}的 $L$-挠阿贝尔\Term{sheaf}{层} $F$，\Term{specialization}{特殊化}\Term{homomorphism}{同态} $(*)$ 在 $q=n+1$ 时是\Term{injectivity}{单射}。
\end{SGAresult}
\begin{proof}
存在\Term{strictly-local}{严格局部}的\Term{integral}{整}\Term{scheme}{概形} $Y'$ 及\Term{morphism}{态射} $g:Y'\to Y$，把\Term{closed-point}{闭点} $y'_0$ 映到 $y_0$、\Term{generic-point}{泛点}映到 $y'_1$：例如先取 $Y$ 关于 $\bar y_0$ 的\Term{strict-localization}{严格局部化}，再取其由 $y_1$ 上方一点定义的\Term{integral}{整}闭子\Term{scheme}{概形}即可。此外，把 $Y'$ 换成它在其\Term{function-field}{函数域}的一个\Term{algebraically-closed}{代数闭}包中的\Term{normalization}{正规化}，可设 $Y'$ \Term{normal}{正规}且 $k(y'_1)$ \Term{algebraically-closed}{代数闭}，从而 $y'_0,y'_1$ 均为 $Y'$ 的\Term{geometric-point}{几何点}。对 $(f,F,g)$ 使用\Term{base-change}{基变换}定理（第 XII 讲 5.1），为证明 2.3，可把 $Y,\bar y_0,\bar y_1$ 换成 $Y',y'_0,y'_1$，从而归结到 $Y$ \Term{normal}{正规}且\Term{strictly-local}{严格局部}、$\bar y_0,\bar y_1$ 分别是其\Term{closed-point}{闭点} $y_0$ 和\Term{generic-point}{泛点} $y_1$ 的情形。另外，利用第 XII 讲 4，可见此时\Term{specialization}{特殊化}\Term{homomorphism}{同态} $(*)$ 等同于\Term{homomorphism}{同态}
\[
H^q(X,F)\longrightarrow H^q(X_1,g_1^*F),
\]
其中 $X_1=X_{y_1}$，$g:y_1\to Y$ 和 $g_1:X_1\to X$ 为\Term{canonical}{典范}\Term{morphism}{态射}。于是断言由下面的引理推出（其中不再假设 $f$ \Term{proper}{固有}）。
\end{proof}
\begin{SGAresult}{引理}{2.4}
沿用 2.3 的预备记号，去掉 $f$ 的\Term{proper}{固有}性假设，改为假设 $Y$ \Term{normal}{正规}\Term{integral}{整}，其\Term{generic-point}{泛点} $y_1$ 的\Term{residue-field}{剩余域} $k(y_1)$ \Term{separably-closed}{可分闭}；令 $X_1=X_{y_1}$，$g_1:X_1\to X$ 为\Term{canonical}{典范}\Term{morphism}{态射}。则\Term{homomorphism}{同态}
\[
H^q(X,F)\longrightarrow H^q(X_1,g_1^*F)\tag{**}
\]
在 $q=0$ 时为\Term{bijection}{双射}（相应地，$q=0,1$ 时为\Term{bijection}{双射}；相应地，$q=n$ 时为\Term{bijection}{双射}，$q=n+1$ 时为\Term{injectivity}{单射}）。此外，当 $n=0$，即 $g$ \Term{locally-acyclic}{局部无环}时，对 $X$ 上任意群\Term{sheaf}{层} $F$，\Term{homomorphism}{同态} $(**)$ 在 $q=1$ 时为\Term{injectivity}{单射}。
\end{SGAresult}
\begin{proof}
一个直接的论证（第 XV 讲 1.6 (i)$\Rightarrow$(ii)）表明，只须证明关系
\[
F\xrightarrow{\sim}g_{1*}g_1^*F,
\]
且在 $n\geq1$ 时还须证明
\[
R^qg_{1*}g_1^*F=0\qquad(1\leq q\leq n).
\]
这些关系对\Term{etale-topology}{平展拓扑}而言在 $X$ 上是\Term{local-scheme}{局部}的，而 $F$ 对\Term{etale-topology}{平展拓扑}\Term{locally-constant}{局部常值}，故可设 $F$ 常值，亦即具有 $G_X$ 的\Term{form}{形式}，其中 $G$ 为\Term{set}{集合}（相应地，$\operatorname{ind}$-$L$-群；相应地，$L$-挠\Term{abelian-group}{阿贝尔群}）。另一方面，$X$ \Term{normal}{正规}的假设蕴含 $g_*(G_{y_1})\xleftarrow{\sim}G_Y$（第 IX 讲 2.14.1），
\SourcePage{148}
而 $k(y_1)$ \Term{separably-closed}{可分闭}的假设蕴含 $R^qg_*(G_{y_1})=0$（$q\geq1$）。另外，$g$ \Term{locally-acyclic}{局部 $0$-无环}（相应地，……）的假设蕴含 $R^qg_*(G_{y_1})$ 的形成在 $q\leq n$ 时与\Term{base-change}{基变换} $f:X\to Y$ 交换；这里用到可设 $Y$ \Term{affine}{仿射}，而此时 $y_1$ 是其\Term{affine}{仿射}\Term{open-neighborhood}{开邻域}的\Term{inverse-limit}{投射极限}，这些\Term{neighborhood}{邻域}在 $Y$ 中\Term{quasi-compact}{拟紧}，因而可应用第 XV 讲 1.11 的定义及取\Term{limit}{极限}理论（第 VII 讲 5）。由此确得
\[
G_X\xrightarrow{\sim}g_{1*}G_{X_1},\qquad R^qg_{1*}G_{X_1}=0\quad(1\leq q\leq n),
\]
故完成 2.4、从而也完成 2.1 的证明。
\end{proof}
\textbf{注记。}上述通过 2.4 的证明比我们最初的证明简单，它出自 M. Lubkin。
\begin{SGAresult}{推论}{2.5}
在 2.1 的条件下，设 $Y$ \Term{connected}{连通}，且 $q\leq n$。则当 $\bar y$ 遍历 $Y$ 的\Term{geometric-point}{几何点}时，各 $H^q(X_{\bar y},F)$ 两两\Term{isomorphism}{同构}。
\end{SGAresult}
\SGAsection{3}{相对上同调纯性定理}
\begin{SGAresult}{定义}{3.1}
设 $S$ 为\Term{scheme}{概形}。所谓 $S$-\Term{smooth-pair}{光滑对} $(Y,X)$，是指\Term{smooth}{光滑} $S$-\Term{prescheme}{预概形}之间的一个 $S$-\Term{closed-immersion}{闭浸入} $i:Y\to X$，亦即\Term{prescheme}{预概形}的\Term{commutative-diagram}{交换图}
\[
\begin{tikzcd}
Y\arrow[r,"i"]\arrow[dr,"h"']&X\arrow[d,"f"]\\&S
\end{tikzcd}
\]
其中 $f,h$ \Term{smooth}{光滑}，$i$ 为\Term{closed-immersion}{闭浸入}。记 $U=X-Y$，$j:U\to X$ 为\Term{open-immersion}{开浸入}，$g:U\to S$ 为\Term{structural-morphism}{结构态射}。称\Term{fiber}{纤维} $Y_s$ 在 $X_s$ 中于 $y\in Y$ 处的\Term{codimension}{余维数}为 $(Y,X)$ 在 $y$ 处的\Term{codimension}{余维数}，其中 $s=h(y)$。
\end{SGAresult}
因此，\Term{codimension}{余维数}是 $Y$ 上的\Term{locally-constant}{局部常值}函数。说“$(Y,X)$ 的\Term{codimension}{余维数}为 $c$”，意指该函数事实上为常数 $c$。

设 $(Y,X)$ 为 $S$-\Term{smooth-pair}{光滑对}，$y\in Y$。回顾（SGA II 4.10），存在整数 $m,n$、$y$ 在 $X$ 中的\Term{neighborhood}{邻域} $X'$ 以及\Term{etale}{平展}\Term{morphism}{态射}
\[
\phi:X'\longrightarrow E_S^n=\operatorname{Spec}\mathcal O_S[t_1,\ldots,t_n],
\]
使 $Y'=Y\cap X'$ 是由方程
\[
t_{m+1}=\cdots=t_n=0
\]
定义的闭子\Term{scheme}{概形} $E_S^m$ 的\Term{inverse-image}{逆像}。
\SourcePage{149}
换言之，在\Term{etale-topology}{平展拓扑}下\Term{local-scheme}{局部}地，每个 $S$-\Term{smooth-pair}{光滑对} $(Y,X)$ 都与适当的 $m,n$ 对应的标准对 $(E_S^m,E_S^n)$ \Term{isomorphism}{同构}，其中
\[
E_S^r=\operatorname{Spec}\mathcal O_S[t_1,\ldots,t_r].
\]
本节要对 $S$-\Term{smooth-pair}{光滑对} $(Y,X)$，计算以\Term{locally-constant}{局部常值}且与\Term{residue-characteristic}{剩余特征}互素的挠阿贝尔\Term{sheaf}{层} $F$ 为系数的\Term{local-cohomology}{局部上同调}\Term{sheaf}{层} $\mathcal H_Y^q(X,F)=(R^qi^!)F$（第 V 讲 4 及第 VIII 讲 6.6）。先作一些预备讨论。
\begin{SGAresult}{命题}{3.2}
设
\[
\begin{tikzcd}
U\arrow[r,"j"]\arrow[dr,"g"']&X\arrow[d,"f"]\\&S
\end{tikzcd}
\]
为\Term{commutative-diagram}{交换图}，$f$ \Term{smooth}{光滑}，$j$ 为\Term{open-immersion}{开浸入}，且对每个 $s\in S$，\Term{fiber}{纤维} $U_s$ 在 $X_s$ 中\Term{dense}{稠密}。设 $F$ 为 $S$ 上的\Term{set}{集合}\Term{sheaf}{层}。则\Term{canonical}{典范}\Term{morphism}{态射}
\[
f^*F\longrightarrow j_*g^*F=j_*j^*f^*F
\]
为\Term{bijection}{双射}。
\end{SGAresult}
\begin{proof}
容易归结到 $U$ 在 $X$ 中\Term{retrocompact}{回紧}的情形。设 $\bar s$ 为 $S$ 的\Term{geometric-point}{几何点}，$\bar x$ 为 $X$ 在 $\bar s$ 上方的\Term{geometric-point}{几何点}，$\widetilde f:\widetilde X\to\widetilde S$ 为 $S,X$ 在相应\Term{geometric-point}{几何点}处的\Term{strict-localization}{严格局部化}间的\Term{morphism}{态射}，且 $\widetilde U=U\times_X\widetilde X$。使用显然的记号，有
\[
(f^*F)_{\bar x}=H^0(\widetilde X,\widetilde f^*\widetilde F)\simeq H^0(\widetilde S,\widetilde F).
\]
故只须证明 $\widetilde g:\widetilde U\to\widetilde S$ 为 $0$-\Term{acyclic}{无环}的，这将蕴含
\[
(j_*g^*F)_{\bar x}\simeq H^0(\widetilde U,\widetilde g^*\widetilde F)\simeq H^0(\widetilde S,\widetilde F).
\]
$\widetilde g$ 是\Term{smooth}{光滑}\Term{morphism}{态射}的\Term{limit}{极限}，故\Term{locally-acyclic}{局部 $0$-无环}（第 XV 讲 2.1、1.11）；由第 XV 讲 1.12，只须证明 $\widetilde U/\widetilde S$ 的\Term{geometric-fiber}{几何纤维} $0$-\Term{acyclic}{无环}，即\Term{connected}{连通}且\Term{nonempty}{非空}。设 $\bar s'$ 为 $\widetilde S$ 的\Term{geometric-point}{几何点}。\Term{fiber}{纤维} $\widetilde X_{\bar s'}$ \Term{regular}{正则}、\Term{connected}{连通}且\Term{nonempty}{非空}；由假设，$\widetilde U_{\bar s'}$ 是 $\widetilde X_{\bar s'}$ 的\Term{dense-open-set}{稠密开集}，因此\Term{connected}{连通}且\Term{nonempty}{非空}，遂得结论。
\end{proof}
\begin{SGAresult}{推论}{3.2.1}
在 3.2 的条件下，由 $X$ 的\Term{etale-covering}{平展覆盖}\Term{category}{范畴}到 $U$ 的\Term{etale-covering}{平展覆盖}\Term{category}{范畴}的\Term{functor}{函子} $X'\mapsto X'|U$ 为\Term{fully-faithful}{全忠实}。特别地，对任意\Term{finite-group}{有限群} $G$，\Term{restriction}{限制}映射
\[
H^1(X,G)\longrightarrow H^1(U,G)
\]
是\Term{injectivity}{单射}。
\end{SGAresult}
\SourcePage{150}
这由 3.2 经一个标准论证而得。为便于引用，将它明确写成一个引理。
\begin{SGAresult}{引理}{3.2.2}
设 $f:X\to Y$ 为\Term{scheme}{概形}\Term{morphism}{态射}，$f^*:\operatorname{Et}(Y)\to\operatorname{Et}(X)$ 为\Term{inverse-image-functor}{逆像函子} $Y'\mapsto Y'\times_YX$，从 $Y$ 的\Term{etale-covering}{平展覆盖}\Term{category}{范畴} $\operatorname{Et}(Y)$ 到 $\operatorname{Et}(X)$。下列条件\Term{equivalence}{等价}：

(i) $f^*$ 忠实（相应地，\Term{fully-faithful}{全忠实}）。

(ii) 对 $Y$ 的任意\Term{etale-covering}{平展覆盖} $Y'$，以 $f':X'\to Y'$ 记由 $f$ 经\Term{base-change}{基变换} $Y'\to Y$ 得到的\Term{morphism}{态射}，则由 $Y'$ 的开\Term{closed-subset}{闭子集}全体到 $X'$ 的开\Term{closed-subset}{闭子集}全体的映射 $U\mapsto f'^{-1}(U)$ 是\Term{injectivity}{单射}（相应地，\Term{bijection}{双射}）。

(ii bis) 沿用 (ii) 的记号，对 $Y'$ 上任意常值\Term{set}{集合}\Term{sheaf}{层} $C$，\Term{canonical}{典范}映射
\[
H^0(Y',C)\longrightarrow H^0(X',f^*C)
\]
是\Term{injectivity}{单射}（相应地，\Term{bijection}{双射}）。

(iii) 沿用 (ii) 的记号，\Term{canonical}{典范}映射
\[
\Gamma(Y'/Y)\longrightarrow\Gamma(X'/X)
\]
是\Term{injectivity}{单射}（相应地，\Term{bijection}{双射}）。
\end{SGAresult}
证明留给读者（参见 SGA 2 第 IX 讲 3.1 和 3.2）。

\textbf{注记 3.2.3。}可以大幅推广 3.2：既可在\Term{set}{集合}\Term{sheaf}{层}的语言中弱化对 $f$ 的假设，也可在群\Term{sheaf}{层}的语言中给出关于 $R^qg_*(f^*F)$ 的结论（参见 SGA 2 第 XIV 讲 1 及 SGA 1 第 XIII 讲）。对 3.4 至 3.6 也有类似注记。

再“回顾” Zariski--Nagata \Term{purity}{纯性}定理的相对\Term{form}{形式}：
\begin{SGAresult}{定理}{3.3}
（相对\Term{purity-theorem}{纯性定理}。）设
\[
\begin{tikzcd}
U\arrow[r,"j"]\arrow[dr,"g"']&X\arrow[d,"f"]\\&S
\end{tikzcd}
\]
为\Term{commutative-diagram}{交换图}，$f$ \Term{smooth}{光滑}，$j$ 为\Term{open-immersion}{开浸入}，且对一切 $s\in S$，$X_s-U_s$ 在 $X_s$ 中处处\Term{codimension}{余维数} $\geq2$。则从 $X$ 的\Term{etale-covering}{平展覆盖}\Term{category}{范畴}到 $U$ 的\Term{etale-covering}{平展覆盖}\Term{category}{范畴}的\Term{functor}{函子} $X'\mapsto X'|U$ 是\Term{category-equivalence}{范畴等价}。
\end{SGAresult}
\SourcePage{151}
\begin{proof}
\Term{functor}{函子}的\Term{fully-faithful}{全忠实}性已见于 3.2.1。\Term{essentially-surjective}{本质满射}性将由经典定理推出。给定 $U$ 的\Term{etale-covering}{平展覆盖} $U'$，须证明它可延拓为 $X$ 的\Term{etale-covering}{平展覆盖} $X'$。容易立即归结到 $U$ 在 $X$ 中\Term{retrocompact}{回紧}的情形，再归结到 $S$ \Term{affine}{仿射}且在 $\operatorname{Spec}\mathbb Z$ 上\Term{finite-type}{有限型}、$X$ 是\Term{connected}{连通}\Term{scheme}{概形}、$f$ \Term{finite-type}{有限型}的情形。

对 $n=\dim S$ 归纳；当 $n=0$ 时，断言由经典定理（SGA 2 第 X 讲 3.4）而得。已有唯一性，故可对\Term{normalization}{正规化} $\overline S\to S$ 应用 SGA 第 IX 讲 4.7 的\Term{descent}{下降}技术，将问题归结到 $S$（因而 $X$）\Term{normal}{正规}的情形（$\overline S\to S$ \Term{finite}{有限}，见 EGA IV 7.8.3 (iii)、(vi)）。此外，适当扩大 $U$，可设 $U$ 在 $U'$ 可延拓的\Term{open-subset}{开集}之中为极大；于是须证明 $U=X$。否则，设 $x$ 为 $X-U$ 的一个\Term{maximal-point}{极大点}，$s=f(x)$。只须证明可把 $U'$ 延拓到 $x$ 的一个\Term{neighborhood}{邻域}上。可设 $\dim\mathcal O_{S,s}=n\geq1$，并且定理对\Term{dimension}{维数} $<n$ 已证明。令 $B=\mathcal O_{X,x}$，$B'$ 为 $B$ 在 $U$ 的\Term{covering}{覆盖} $U'$ 上的\Term{rational-function-ring}{有理函数环}中的\Term{normalization}{正规化}；则 $B'$ 是\Term{finite}{有限} $B$-\Term{algebra}{代数}（同上）。显然只须证明 $B'$ 是\Term{etale}{平展} $B$-\Term{algebra}{代数}。

设 $A=\mathcal O_{S,s}$，$0\ne t\in\operatorname{max}A$，$A_0=A/tA$，$B_0=B/tB$。由归纳假设，\Term{covering}{覆盖}在 $U_0=U\otimes_AA_0$ 上的\Term{restriction}{限制}可延拓为 $\operatorname{Spec}B_0$ 的\Term{etale-covering}{平展覆盖}，记作 $\operatorname{Spec}B_0^*$。允许将 $X$ 换成 $x$ 的一个\Term{neighborhood}{邻域}的适当\Term{etale-covering}{平展覆盖}，因而可设 $\operatorname{Spec}B_0^*$ 在 $\operatorname{Spec}B_0$ 上\Term{completely-split}{完全分裂}。于是归结到 $B_0$-\Term{algebra}{代数} $B'_0$ 在 $\operatorname{Spec}B_0$ 的\Term{closed-point}{闭点}之外\Term{completely-split}{完全分裂}的情形。这时第 XV 讲 3.3 的假设满足，遂得结论。
\end{proof}
\begin{SGAresult}{推论}{3.4}
假设同 3.3。则对 $S$ 上任意\Term{ind-finite-group}{ind-有限}群\Term{sheaf}{层} $F$，有 $(R^1j_*)g^*F=0$。
\end{SGAresult}
\begin{proof}
如常归结到 $S$ \Term{noetherian}{诺特}、$f$ \Term{finite-type}{有限型}、$F$ \Term{constructible}{可构造}的情形。设 $0\to F\to G$ 为 $S$ 上\Term{ind-finite-group}{ind-有限}群\Term{sheaf}{层}的\Term{injectivity}{单射}，$C=G/F$。由 3.2，\Term{functor}{函子} $j_*g^*$ \Term{exact}{正合}，故 $j_*g^*G\to j_*g^*C$ \Term{surjectivity}{满射}，从而（第 XII 讲 3.1）$(R^1j_*)g^*F\to(R^1j_*)g^*G$ \Term{injectivity}{单射}。因此可把 $F$ 换成 $G$，由第 IX 讲 2.14 及第 VIII 讲 5.5，归结到 $F=G$ 是取值于普通\Term{finite-group}{有限群} $G$ 的\Term{constant-sheaf}{常值层}的情形。设 $\bar x$ 为 $X$ 的\Term{geometric-point}{几何点}，$\widetilde X$ 为 $X$ 在 $\bar x$ 处的\Term{strict-localization}{严格局部化}，$\widetilde U=\widetilde X\times_XU$。由 3.3 有
\[
(R^qj_*g^*F)_{\bar x}=H^q(\widetilde U,G_{\widetilde U})\simeq H^q(\widetilde X,G_{\widetilde X}),
\]
而最后一项为零，遂得结论。
\end{proof}
\begin{SGAresult}{引理}{3.5}
（相对 Abhyankar 引理。）设 $(Y,X)$ 为\Term{codimension}{余维数} $1$ 的 $S$-\Term{smooth-pair}{光滑对}，$Y$ 在 $X$ 中由一个方程 $t=0$ 定义。设 $U=X-Y$，
\SourcePage{152}
$V/U$ 为次数 $n$ 与 $Y$ 的\Term{residue-characteristic}{剩余特征}互素的主 Galois \Term{covering}{覆盖}。令
\[
X'=\operatorname{Spec}\bigl(\mathcal O_X[z]/(z^n-t)\bigr),\quad U'=U\times_XX',\quad V'=V\times_XX'.
\]
则\Term{covering}{覆盖} $V'$ 唯一延拓为 $X'$ 的一个\Term{etale-covering}{平展覆盖}。
\end{SGAresult}
\begin{proof}
立即归结到 $S,X,Y$ 均\Term{affine}{仿射}且在 $\operatorname{Spec}\mathbb Z$ 上\Term{finite-type}{有限型}的情形。唯一性是 3.2.1 的特例。由第 XV 讲 1.12 推出存在性：设 $\bar s$ 为 $S$ 的\Term{geometric-point}{几何点}，$\bar x$ 为 $Y$ 在 $\bar s$ 上方的\Term{geometric-point}{几何点}。设 $\widetilde f:\widetilde X\to\widetilde S$ 为 $f$ 诱导的相应\Term{strict-localization}{严格局部化}间的\Term{morphism}{态射}，$\widetilde U=U\times_X\widetilde X$，其余类推。考虑到唯一性，由\Term{descent}{下降}，显然只须证明 $\widetilde V'$ 可延拓为 $\widetilde X'$ 的\Term{etale-covering}{平展覆盖}；由于 $\widetilde X'$ 没有非\Term{trivial-covering}{平凡覆盖}，这就是证明 $\widetilde V'$ 为 $\widetilde U'$ 的\Term{trivial-covering}{平凡覆盖}。

令
\[
\overline U=\varprojlim_m\widetilde U_m,\qquad \widetilde U_m=\operatorname{Spec}\mathcal O_{\widetilde U}[t^{1/m}],
\]
其中 $m$ 遍历与 $\widetilde X$ 的\Term{residue-characteristic}{剩余特征}互素的整数。则 $\overline U/\widetilde U$ 是具有\Term{abelian-group}{阿贝尔群} $G=\varprojlim_m\mathbb Z/m\mathbb Z$ 的“无穷 \Term{galois-covering}{Galois 覆盖}”，立即可见，只须证明 $\overline U$ 对 $\widetilde X$ 的\Term{residue-characteristic}{剩余特征}之补集 $L\subset\mathbb P$ 是 $1$-无球的。

考虑\Term{morphism}{态射} $\overline f:\overline U\to\widetilde S$。它是\Term{smooth}{光滑}\Term{morphism}{态射}的\Term{limit}{极限}，因而对 $L$ \Term{locally-aspherical}{局部 $1$-无球}（第 XV 讲 2.1、1.11 (ii)）。由第 XV 讲 1.12，只须证明 $\overline U/\widetilde S$ 的\Term{geometric-fiber}{几何纤维}对 $L$ $1$-无球。设 $\bar s$ 为 $\widetilde S$ 的\Term{geometric-point}{几何点}。\Term{geometric-fiber}{几何纤维} $\overline U_{\bar s}$ 是各 $(\widetilde U_m)_{\bar s}$ 的\Term{limit}{极限}。因此，只须证明每个 $(\widetilde U_m)_{\bar s}$ 上次数与\Term{residue-characteristic}{剩余特征}互素的\Term{principal-galois-covering}{主 Galois 覆盖}，都在 $\overline U_{\bar s}$ 上诱导\Term{trivial-covering}{平凡覆盖}。把 $\widetilde X$ 换成 $\widetilde X_m=\operatorname{Spec}\mathcal O_{\widetilde X}[t^{1/m}]$（它在 $\widetilde S$ 上仍\Term{smooth}{光滑}），归结到 $m=1$，即 $\widetilde U_m=\widetilde U$ 的情形。设 $V$ 为 $\widetilde U_{\bar s}$ 的一个这种\Term{covering}{覆盖}。$\widetilde X_{\bar s}=Z$ 为\Term{strictly-local}{严格局部}的\Term{regular}{正则}\Term{scheme}{概形}，而 $W=\widetilde U_{\bar s}$ 是一个\Term{codimension}{余维数} $1$ 的\Term{regular}{正则}\Term{subset}{子集}的补\Term{open-subset}{开集}。因此可应用通常\Term{form}{形式}的 Abhyankar 引理（SGA 1 第 X 讲 3.6），得对适当的 $m$，$V$ 在 $W_m$ 上诱导的\Term{covering}{覆盖} $V_m$ 可延拓为 $Z_m$ 的\Term{etale-covering}{平展覆盖}。但 $Z$、从而 $Z_m$ \Term{strictly-local}{严格局部}，故该\Term{covering}{覆盖}平凡，因而 $V_m$ 是 $W_m$ 的\Term{trivial-covering}{平凡覆盖}。
\end{proof}
\textbf{注记 3.5.1。}若不假设 $G$ 的阶 $n$ 与 $X$ 的\Term{residue-characteristic}{剩余特征}互素，仍可证明：令 $m$ 为整除 $G$ 的阶且与 $Y$ 的\Term{residue-characteristic}{剩余特征}互素的最大整数，并假设对每个点 $s\in S$，$U_s=X_s-Y_s$ 的\Term{covering}{覆盖} $V_s$ 在 $Y_s$ 的\Term{maximal-point}{极大点}处的\Term{inertia-group}{惯性群}的阶均与 $Y$ 的\Term{residue-characteristic}{剩余特征}互素。则取 $X'=\operatorname{Spec}(\mathcal O_X[z]/(z^m-t))$，3.5 的结论仍成立。更一般地，还可取 $m$ 为刚才那些\Term{inertia-group}{惯性群}的阶的任意公倍数，这些阶假设与相应的\Term{residue-characteristic}{剩余特征}互素。
\SourcePage{153}
\begin{SGAresult}{推论}{3.6}
（\Term{dimension}{维数} $1$ 的\Term{cohomology}{上同调}\Term{purity}{纯性}。）设 $(Y,X)$ 为\Term{codimension}{余维数} $1$ 的 $S$-\Term{smooth-pair}{光滑对}，$F$ 为 $X$ 上\Term{constructible}{可构造}且\Term{locally-constant}{局部常值}的群\Term{sheaf}{层}，其阶与 $X$ 的\Term{residue-characteristic}{剩余特征}互素。则沿用 3.1 的记号，$(R^1j_*)j^*F$ 作为\Term{pointed-set}{带基点集合}\Term{sheaf}{层}，\Term{local-isomorphism}{局部同构}于 $i_*(F_Y/\operatorname{int}(F_Y))$，其中 $F_Y=F|Y$，$F_Y/\operatorname{int}(F_Y)$ 表示 $F_Y$ 对其自身以\Term{inner-automorphism}{内自同构}\Term{action}{作用}所得之商。
\end{SGAresult}
\textbf{注记。}推论 3.6 的\Term{isomorphism}{同构}不是\Term{canonical}{典范}的。
\begin{proof}
立即归结到 $X$ \Term{noetherian}{诺特}且\Term{connected}{连通}的情形。断言对\Term{etale-topology}{平展拓扑}而言在 $X$ 上是\Term{local-scheme}{局部}的，故可设 $F$ 为取值于普通\Term{finite-group}{有限群} $G$ 的\Term{constant-sheaf}{常值层}，$G$ 的阶 $m$ 与 $X$ 的\Term{residue-characteristic}{剩余特征}互素，并可进一步设 $\mu_{mX}\simeq(\mathbb Z/m\mathbb Z)_X$。令 $U=X-Y$，$U'=\operatorname{Spec}\mathcal O_U[t^{1/m}]$；由于 $\mu_{mX}\simeq(\mathbb Z/m\mathbb Z)_X$，它是 $U$ 的一个群为 $\mathbb Z/m\mathbb Z$ 的“本质”\Term{principal-galois-covering}{主 Galois 覆盖}。$U$ 上那些被 $U'$ \Term{trivialization}{平凡化}的群 $G$ 的\Term{principal-covering}{主覆盖}，在 $Y$ 上\Term{local-scheme}{局部}地由
\[
H=\operatorname{Hom}(\mathbb Z/m\mathbb Z,G)\bmod\operatorname{int}(G)
\]
分类；它是与 $G/\operatorname{int}(G)$ 的\Term{underlying-set}{底集}合相\Term{isomorphism}{同构}的\Term{pointed-set}{带基点集合}。由此得到从 $Y$ 上\Term{constant-sheaf}{常值层} $H_Y$ 到 $(R^1j_*)G_U$ 的\Term{morphism}{态射}，只须证明它为\Term{bijection}{双射}，而这立即由 3.5 而得。
\end{proof}
\begin{SGAresult}{定理}{3.7}
（\Term{cohomology}{上同调}\Term{purity}{纯性}。）设 $(Y,X)$ 为\Term{codimension}{余维数} $c>0$ 的 $S$-\Term{smooth-pair}{光滑对}。设 $F$ 为 $X$ 上的阿贝尔\Term{sheaf}{层}，在\Term{etale-topology}{平展拓扑}下\Term{local-isomorphism}{局部同构}于 $f^*G$ \Term{form}{形式}的\Term{sheaf}{层}，其中 $G$ 为 $S$ 上与\Term{residue-characteristic}{剩余特征}互素的挠\Term{sheaf}{层}（例如 $F$ 是 $X$ 上\Term{locally-constant}{局部常值}的、阶与 $X$ 的\Term{residue-characteristic}{剩余特征}互素的\Term{finite-group}{有限群}\Term{sheaf}{层}）。则沿用 3.1 的记号，有
\[
\mathcal H_Y^q(X,F)=(R^qi^!)F=0\qquad(q\ne2c),
\]
亦即
\[
F\xrightarrow{\sim}j_*j^*F,\qquad (R^qj_*)j^*F=0\quad(q\ne0,2c-1),
\]
并且
\[
\mathcal H_Y^{2c}(X,F)=i^*[(R^{2c-1}j_*)j^*F]
\]
是\Term{local-isomorphism}{局部同构}于 $i^*F$ 的\Term{sheaf}{层}。
\end{SGAresult}
\begin{proof}
断言对\Term{etale-topology}{平展拓扑}而言在 $X$ 上是\Term{local-scheme}{局部}的；故可设 $F$ 常值\footnote{至少当假设 $F$ \Term{locally-constant}{局部常值}时如此。读者会确信，下面的证明本质上也适用于一般情形。}，因而如愿还可设 $F\simeq f^*(\mathbb Z/n\mathbb Z)_S$，其中 $(\mathbb Z/n\mathbb Z)_S$ 是 $S$ 上的\Term{constant-sheaf}{常值层}。

先处理 $c=1$。（关于 $(R^1j_*)j^*F$ 的断言，至少在 $F$ \Term{locally-constant}{局部常值}时，可立即由 3.6 得到，但我们将重新推导它。）
\SourcePage{154}
回顾，$(Y,X)$ \Term{local-isomorphism}{局部同构}于对 $(E_S^{n-1},E_S^n)$，把 $S$ 换成 $E_S^{n-1}$，即归结到 $X=E_S^1=\operatorname{Spec}\mathcal O_S[t]$、$Y$ 为 $X/S$ 的\Term{section}{截面} $t=0$ 的情形。这一情形\Term{local-isomorphism}{局部同构}于\Term{projective-space}{射影空间} $\mathbb P_S^1$ 的无穷远\Term{section}{截面}的\Term{embedding}{嵌入}。因此可取 $X=\mathbb P_S^1$，$Y$ 为无穷远\Term{section}{截面}，$U=E_S^1=X-Y$。

考察 Leray \Term{spectral-sequence}{谱序列}
\[
(R^pf_*)(R^qj_*)(\mathbb Z/n)\Longrightarrow(R^{p+q}g_*)(\mathbb Z/n).
\]
由第 XV 讲 2.2，当 $p+q>0$ 时其\Term{abutment}{终项}为零。而当 $q>0$ 时，各 $R^qj_*$ 集中在 $Y$ 上，$Y$ 经 $f$ \Term{isomorphism}{同构}地映到 $S$；另一方面，由 3.2 有 $j_*(F|U)\xleftarrow{\sim}F$。所以当 $p,q>0$ 时 $(R^pf_*)(R^qj_*)=0$；当 $q>0$ 时 $f_*(R^qj_*)$ 决定\Term{sheaf}{层} $R^qj_*$；最后 $R^pf_*j_*(F|U)=R^pf_*F$。故\Term{spectral-sequence}{谱序列}化为\Term{isomorphism}{同构}
\[
f_*(R^qj_*)(F|U)\xrightarrow{\sim}(R^{q+1}f_*)F\qquad(q>0).
\]
但 $(R^qf_*)F$ 可逐\Term{fiber}{纤维}计算（第 XII 讲 5.2），因而当 $q\ne0,2$ 时为零，当 $q=2$ 时是\Term{local-isomorphism}{局部同构}于 $G$ 的\Term{sheaf}{层}，这容易由第 XII 讲 5.2 和第 IX 讲 4.7 得出。因此当 $q\ne1$ 时 $(R^qj_*)(F|U)=0$，当 $q=1$ 时它是\Term{local-isomorphism}{局部同构}于 $i^*F$ 的\Term{sheaf}{层}。

现在可用归纳法容易地处理 $c>1$ 的情形：设 $(Y,X)$ 为\Term{codimension}{余维数} $c>1$ 的 $S$-\Term{smooth-pair}{光滑对}。显然，在 $X$ 上\Term{local-scheme}{局部}地可找到一个包含 $Y$、在 $S$ 上\Term{smooth}{光滑}的\Term{codimension}{余维数} $1$ 的子\Term{scheme}{概形} $Z$，使有\Term{commutative-diagram}{交换图}
\[
\begin{tikzcd}
Y\arrow[r,"u"]\arrow[dr,"i"']&Z\arrow[d,"v"]\\&X
\end{tikzcd}
\]
其中每一对都是 $S$-\Term{smooth-pair}{光滑对}（应称 $(Y,Z,X)$ 为 $S$-\Term{smooth}{光滑}三元组），$(Z,X)$ 的\Term{codimension}{余维数}为 $1$，$(Y,Z)$ 的\Term{codimension}{余维数}为 $c-1$。

有\Term{composite-functor}{复合函子}的\Term{spectral-sequence}{谱序列}
\[
(R^pu^!)(R^qv^!)F\Longrightarrow(R^{p+q}i^!)F.
\]
由归纳假设，$(R^qv^!)F$ 在 $q\ne2$ 时为零，在 $q=2$ 时\Term{local-isomorphism}{局部同构}于 $v^*F$；再对 $(R^qv^!)F$ 应用归纳假设，得 $(R^pu^!)(R^qv^!)F$ 当 $(p,q)\ne(2c-2,2)$ 时为零，当 $(p,q)=(2c-2,2)$ 时\Term{local-isomorphism}{局部同构}于 $u^*w^*F$，立即得出关于 $(R^{p+q}i^!)F$ 的结论。
\end{proof}
\begin{SGAresult}{推论}{3.8}
设 $(Y,X)$ 为\Term{codimension}{余维数} $c$ 的 $S$-\Term{smooth-pair}{光滑对}，$n$ 为与 $X$ 的\Term{residue-characteristic}{剩余特征}互素的自然数，令
\[
T_{Y/X}=\mathcal H_Y^{2c}(\mathbb Z/n\mathbb Z).
\]
\SourcePage{155}
由 (3.7)，它是 $Y$ 上在\Term{etale-topology}{平展拓扑}下\Term{local-isomorphism}{局部同构}于 $(\mathbb Z/n\mathbb Z)_Y$ 的\Term{sheaf}{层}。设 $F$ 为 $X$ 上被 $n$ 零化且满足 3.7 中条件的阿贝尔\Term{sheaf}{层}。则（重申一下）$\mathcal H_Y^i(F)=0$（$i\ne2c$），并且有\Term{canonical-isomorphism}{典范同构}
\[
\mathcal H_Y^{2c}(F)\simeq i^*F\otimes T_{Y/X}.
\]
\end{SGAresult}
\begin{proof}
利用\Term{canonical-homomorphism}{典范同态} $F\to\mathcal Hom(\mathbb Z/n\mathbb Z,F)$，容易定义\Term{canonical-homomorphism}{典范同态}
\[
\varphi:i^*F\otimes T_{Y/X}\longrightarrow\mathcal H_Y^{2c}(F),
\]
亦即\Term{homomorphism}{同态}
\[
i^*F\longrightarrow\mathcal Hom\bigl(\mathcal H_Y^{2c}(\mathbb Z/n\mathbb Z),\mathcal H_Y^{2c}(F)\bigr).
\]
只须证明 $\varphi$ 为\Term{isomorphism}{同构}；沿着给出的 3.7 的证明即可容易核实。
\end{proof}
\begin{SGAresult}{推论}{3.8.1}
在 3.8 的条件下，$\mathcal H_Y^i(F)$ 的形成与任意\Term{base-change}{基变换} $S'\to S$ 交换。
\end{SGAresult}
这是显然的。
\begin{SGAresult}{推论}{3.9}
（\Term{absolute-cohomological-purity}{绝对上同调纯性}定理。）设 $X$ 为\Term{regular}{正则}\Term{scheme}{概形}，$Y$ 为处处\Term{codimension}{余维数} $c$ 的\Term{regular}{正则}闭子\Term{prescheme}{预概形}。再设 $X$ 在一个\Term{perfect-field}{完美域} $k$ 上\Term{locally-finite-type}{局部有限型}。则若 $n$ 与 $X$ 的\Term{residue-characteristic}{剩余特征}互素，有
\[
\mathcal H_Y^i(\mathbb Z/n\mathbb Z)=0\qquad(i\ne2c),
\]
且 $\mathcal H_Y^{2c}(\mathbb Z/n\mathbb Z)$ 为 $Y$ 上在\Term{etale-topology}{平展拓扑}下\Term{local-isomorphism}{局部同构}于\Term{constant-sheaf}{常值层} $\mathbb Z/n\mathbb Z$ 的\Term{sheaf}{层}。
\end{SGAresult}
这是 3.7 的特例，因为假设蕴含 $X,Y$ 在 $k$ 上\Term{smooth}{光滑}。

\textbf{注记 3.10。}a) 在 3.8 中可以明确写出\Term{sheaf}{层} $T_{Y/X}$ 的结构：它\Term{canonical-isomorphism}{典范同构}于\Term{sheaf}{层} $(\mu_n)_Y^{\otimes(-c)}$，其中 $\mu_n$ 表示 $n$ 次\Term{root-of-unity}{单位根}\Term{sheaf}{层}（\Term{local-isomorphism}{局部同构}于\Term{constant-sheaf}{常值层} $\mathbb Z/n\mathbb Z$）。在研究\Term{duality}{对偶性}的第 XVIII 讲中，将在稍广的相对\Term{complete-intersection}{完全交}情形下，以 $Y$ 在 $X$ 中的\Term{local-fundamental-class}{局部基本类}之名研究这个\Term{isomorphism}{同构}。

b) 即使不要求存在\Term{base-field}{基域} $k$，3.9 和 3.10 a) 的结论也极有可能成立，至少在 $X$ 为“\Term{excellent}{优良}”\Term{prescheme}{预概形}时如此。不论如何，当 $X$ \Term{characteristic}{特征}为零时，第 XIX 讲将利用 Hironaka 的\Term{resolution-of-singularities}{奇点消解}证明这一点。如将在 SGA 5 中看到的，为建立“\Term{local-scheme}{局部}\Term{duality}{对偶}公式”，尤其需要“绝对\Term{purity}{纯性}定理”（以及\Term{resolution-of-singularities}{奇点消解}）；
\SourcePage{156}
该公式本身又是 Lefschetz--Verdier 公式的一个主要组成部分。
\SGAsection{4}{复数域上代数预概形的上同调比较定理}
设 $X$ 为 $\operatorname{Spec}\mathbb C$ 上的\Term{algebra}{代数}\Term{scheme}{概形}。我们将比较 $X$ 上的\Term{etale}{平展}\Term{site}{景}和经典\Term{site}{景}，并沿用第 XI 讲 4 的记号。作为第 XI 讲 4.4 的一般\Term{form}{形式}，\Term{comparison-theorem}{比较定理}如下。
\begin{SGAresult}{定理}{4.1}
设 $f:X\to S$ 为 $\operatorname{Spec}\mathbb C$ 上\Term{locally-finite-type}{局部有限型}\Term{scheme}{概形}之间的\Term{finite-type}{有限型}\Term{morphism}{态射}，因而有\Term{site}{景}\Term{morphism}{态射}的\Term{commutative-diagram}{交换图}
\[
\begin{tikzcd}
X_{\text{ét}}\arrow[d,"f_{\text{ét}}"']&X_{\mathrm{cl}}\arrow[l,"\varepsilon"']\arrow[d,"f_{\mathrm{cl}}"]\\
S_{\text{ét}}&S_{\mathrm{cl}}\arrow[l,"\varepsilon"]
\end{tikzcd}
\]
设 $F$ 为\Term{set}{集合}\Term{sheaf}{层}（相应地，\Term{ind-finite-group}{ind-有限}群\Term{sheaf}{层}；相应地，挠阿贝尔\Term{sheaf}{层}），并假设处于以下情形之一：

(i) $f$ \Term{proper}{固有}；

(ii) $F$ \Term{constructible}{可构造}。

则\Term{base-change}{基变换}\Term{morphism}{态射}（第 XII 讲 4.1.2）\footnote{实际上，这里使用\Term{base-change}{基变换}\Term{homomorphism}{同态}的条件比第 XII 讲所考虑的更一般；该讲只限于\Term{prescheme}{预概形}\Term{morphism}{态射}的\Term{cartesian-square}{笛卡尔方块}。读者容易确信，该定义逐字适用于\Term{site}{景}\Term{morphism}{态射}的本质\Term{commutative-square}{交换方块}。有关进一步的展开，见下一讲。}
\[
\varphi:\varepsilon^*(R^qf_{\text{ét}*})F\longrightarrow(R^qf_{\mathrm{cl}*})\varepsilon^*F
\]
在 $q=0$ 时（相应地，$q=0,1$ 时；相应地，$q\geq0$ 时）为\Term{bijection}{双射}。
\end{SGAresult}
注意，在第 XI 讲 4.4 的证明中所使用的 Grauert--Remmert 定理（第 XI 讲 4.3 (iii)）是 4.1 的特例。

证明时先处理 $f$ \Term{proper}{固有}的情形。此时，结果将由 GAGA 和\Term{proper}{固有}\Term{morphism}{态射}的\Term{base-change}{基变换}定理，经过一个本来也可列于第 XIV 讲的初等论证而得。处理一般情形时，还将使用\Term{resolution-of-singularities}{奇点消解} [1] 和 3.7。

显然可设 $Y$、因而也可设 $X$ 在 $\mathbb C$ 上\Term{finite-type}{有限型}。

\textbf{$f$ \Term{proper}{固有}的情形。}回顾以下初等事实：
\SourcePage{157}
(a) \Term{topological-space}{拓扑空间} $X_{\mathrm{cl}}$ \Term{local-scheme}{局部}紧。

(b) 若 $f:X\to S$ \Term{proper}{固有}，则 $f_{\mathrm{cl}}:X_{\mathrm{cl}}\to S_{\mathrm{cl}}$ 是\Term{topological-space}{拓扑空间}间的\Term{proper-map}{固有映射}。

事实上，(a) 是平凡的，因为\Term{local-scheme}{局部}地 $X_{\mathrm{cl}}$ 是某个 $\mathbb C^n$ 中的闭集。对 (b)，由 Chow 引理归结到 $f$ \Term{projective}{射影}的情形，而此时结论平凡。

设 $F$ 为 $X_{\text{ét}}$ 上的\Term{set}{集合}\Term{sheaf}{层}（相应地，……）。对两个\Term{morphism}{态射} $f_{\text{ét}}$ 和 $f_{\mathrm{cl}}$，$R^q$ 的形成都与取\Term{fiber}{纤维}交换（第 XII 讲 5.2 及 [2] 4.11.1）。所以只须在 $S_{\text{ét}}$ 的\Term{closed-point}{闭点}处逐\Term{fiber}{纤维}核实定理。

在下面的引理中，不假设 $g$ \Term{proper}{固有}。
\begin{SGAresult}{引理}{4.2}
设
\[
\begin{tikzcd}
X\arrow[dr,"g"']&\overline X\arrow[l,"h"']\arrow[d,"\overline g"]\\&S
\end{tikzcd}
\]
为\Term{commutative-diagram}{交换图}，$h$ \Term{proper}{固有}，$U$ 为 $X$ 的\Term{open-subset}{开集}，使 $h$ 将 $\overline X$ 的\Term{open-subset}{开集} $\overline U=U\times_X\overline X$ \Term{isomorphism}{同构}地映到 $U$。设 $j:U\to X$ 和 $\overline j:\overline U\to\overline X$ 为包含，$F$ 为 $U$ 上的\Term{set}{集合}\Term{sheaf}{层}（相应地，……），$\overline F$ 为其在 $\overline U$ 上的诱导\Term{sheaf}{层}。若定理 4.1 对 $(\overline g,\overline j_!\overline F)$ 成立，则也对 $(g,j_!F)$ 成立。
\end{SGAresult}
\begin{proof}
定理对 $(h,\overline j_!\overline F)$ 成立，因为可归结到逐\Term{fiber}{纤维}核实，而此时只有两种情形：或\Term{sheaf}{层}在\Term{fiber}{纤维}上的\Term{restriction}{限制}为零，或\Term{morphism}{态射} $h$ 为\Term{isomorphism}{同构}。这两种情形的结果都是平凡的。于是得到
\[
\varepsilon^*j_!F\simeq\varepsilon^*h_{\text{ét}*}(\overline j_!\overline F)\simeq h_{\mathrm{cl}*}\varepsilon^*(\overline j_!\overline F)
\]
以及对所考虑的 $q$ 值，
\[
(R^qh_{\text{ét}*})\overline j_!\overline F=(R^qf_{\mathrm{cl}*})\varepsilon^*(\overline j_!\overline F)=0\quad(q>0),
\]
故
\[
(R^qg_{\text{ét}*})j_!F\simeq(R^q\overline g_{\text{ét}*})\overline j_!\overline F
\]
及
\[
(R^qg_{\mathrm{cl}*})j_!F\simeq(R^q\overline g_{\mathrm{cl}*})\overline j_!\overline F,
\]
从而得 4.2。
\end{proof}
\begin{SGAresult}{引理}{4.3}
为对所有 $f$ \Term{proper}{固有}的数据 $(f,F)$ 核实 4.1，只须在进一步假设 $f$ 的\Term{relative-dimension}{相对维数} $\leq1$ 且 $S=\operatorname{Spec}\mathbb C$ 为一点的情形下核实它。
\end{SGAresult}
\SourcePage{158}
\begin{proof}
对\Term{relative-dimension}{相对维数}归纳。设结果在\Term{relative-dimension}{相对维数} $<n$ 时已知，且 $n>1$。如前所述，可以逐\Term{fiber}{纤维}核实，即可设 $S=\operatorname{Spec}\mathbb C$ 为一点，因而 $\dim X\leq n$。还可设 $X$ \Term{reduced}{约化}。取 $X$ 上一个在每个\Term{irreducible-component}{不可约分支}上都非常值的\Term{rational-function}{有理函数} $\phi$，$U$ 为 $X$ 上 $\phi$ 有定义的\Term{dense-open-set}{稠密开集}，$\overline X\subset X\times\mathbb P^1$ 为\Term{morphism}{态射} $\phi:U\to E^1$ 的图的\Term{closure}{闭包}。有\Term{commutative-diagram}{交换图}
\[
\begin{tikzcd}
&\overline X\arrow[dl,"h"']\arrow[dr,"a"]\arrow[dd,"\overline g" description]&\\
X\arrow[dr,"g"']&&\mathbb P^1\arrow[dl,"b"]\\
&Y=\operatorname{Spec}\mathbb C&
\end{tikzcd}
\]
其中 $a,b$ 的\Term{relative-dimension}{相对维数}均 $<n$。由归纳假设，定理对 $(a,\cdot)$ 及 $(b,\cdot)$ 成立，故也对 $(\overline g,\cdot)$ 成立。事实上，$q=0$ 时这是平凡的；若 $F$ 是 $\overline X$ 上的挠阿贝尔\Term{sheaf}{层}，则由 Leray \Term{spectral-sequence}{谱序列}之间的\Term{morphism}{态射}
\[
\begin{tikzcd}[column sep=large]
\varepsilon^*(R^ph_{\text{ét}*})(R^qa_{\text{ét}*})F\arrow[r,Rightarrow]\arrow[d]&\varepsilon^*(R^{p+q}\overline g_{\text{ét}*})F\arrow[d]\\
(R^pb_{\mathrm{cl}*}(R^qa_{\mathrm{cl}*}))\varepsilon^*F\arrow[r,Rightarrow]&(R^{p+q}\overline g_{\mathrm{cl}*})\varepsilon^*F
\end{tikzcd}
\]
而得。最后，若 $F$ 为 $\overline X$ 上的\Term{ind-finite-group}{ind-有限}群\Term{sheaf}{层}，注意 $\varepsilon^*(R^1\overline g_{\text{ét}*})$ 是\Term{effaceable}{可消去}\Term{functor}{函子}，故考虑到 $q=0$ 的结果，只须证明 $(R^1\overline g_{\mathrm{cl}*})\varepsilon^*$ 也\Term{effaceable}{可消去}（参见第 XII 讲 8.2）；这由\Term{exact-sequence}{正合列}
\[
0\longrightarrow(R^1b_{\mathrm{cl}*})a_{\mathrm{cl}*}\varepsilon^*F\longrightarrow(R^1\overline g_{\mathrm{cl}*})\varepsilon^*F\longrightarrow b_{\mathrm{cl}*}(R^1a_{\mathrm{cl}*})\varepsilon^*F
\]
以及归纳假设所给出的 $(R^1b_{\mathrm{cl}*})\varepsilon^*$ 和 $(R^1a_{\mathrm{cl}*})\varepsilon^*$ 的\Term{effaceable}{可消去}性而得。

现在设 $F$ 为 $X$ 上的\Term{set}{集合}\Term{sheaf}{层}。有 $F\hookrightarrow h_{\text{ét}*}h_{\text{ét}}^*F=G$。由\Term{exact-sequence}{正合列}
\[
F\longrightarrow G\rightrightarrows G\amalg_FG,\tag{4.3.1}
\]
在 $q=0$ 时归结到对 $G$ 证明 4.1，即对具有 $h_{\text{ét}*}\overline F$ \Term{form}{形式}的\Term{sheaf}{层}证明之。但对这种\Term{sheaf}{层}，结果由定理对 $(\overline g,\cdot)$ 和 $(h,\cdot)$ 的成立而得；立即可见 $h$ 的\Term{relative-dimension}{相对维数} $<n$。
\SourcePage{159}
设 $F$ 为 $X$ 上的挠阿贝尔\Term{sheaf}{层}。有\Term{exact-sequence}{正合列}
\[
0\longrightarrow j_!j^*F\longrightarrow F\longrightarrow i_*i^*F\longrightarrow0.\tag{*}
\]
由于 $U$ 在 $X$ 中\Term{dense}{稠密}，有 $\dim Y\leq n$，故 4.1 对 $(g|Y,i^*F)$ 成立，从而对 $(g,i_*i^*F)$ 成立。因此由\Term{five-lemma}{五引理}，只须对 $(f,j_!j^*F)$ 核实 4.1，而这由定理对 $\overline g$ 的成立推出（4.2）。

设 $F$ 为 $X$ 上的\Term{ind-finite-group}{ind-有限}群\Term{sheaf}{层}。只须证明\Term{functor}{函子} $(R^1g_{\mathrm{cl}*})\varepsilon^*$ 对\Term{sheaf}{层} $F$ \Term{effaceable}{可消去}；这由 $(*)$ 及它对 $j_!j^*F$、$i_*i^*F$ 的\Term{effaceable}{可消去}性而得，即由定理对 $(g,j_!j^*F)$（4.2）及 $(g,i_*i^*F)$（归纳假设）的成立而得。
\end{proof}
\begin{SGAresult}{引理}{4.4}
为在 $f$ \Term{proper}{固有}时证明 (4.1)，只须证明：设 $X$ 为 $\operatorname{Spec}\mathbb C$ 上的\Term{complete}{完备}\Term{regular}{正则}曲线，$F$ 为 $X$ 上常值且\Term{constructible}{可构造}的\Term{set}{集合}\Term{sheaf}{层}（相应地，……），则\Term{morphism}{态射}
\[
H^q(X_{\text{ét}},F)\longrightarrow H^q(X_{\mathrm{cl}},F)
\]
对所考虑的 $q$ 值为\Term{bijection}{双射}。
\end{SGAresult}
\begin{proof}
只须逐\Term{fiber}{纤维}核实定理，故由 4.3 归结到 $X$ \Term{proper}{固有}且\Term{dimension}{维数} $\leq1$ 的情形。\Term{dimension}{维数} $0$ 的情形是平凡的，由此也知定理对\Term{finite-morphism}{有限态射}成立。由于 $f$ \Term{proper}{固有}，$X_{\text{ét}}$ 和 $X_{\mathrm{cl}}$ 的\Term{cohomology}{上同调}与\Term{inductive-limit}{归纳极限}交换（第 VII 讲 3.3 及 [2] 4.12.1），故归结到 $F$ \Term{constructible}{可构造}的情形（第 IX 讲 2.9 (iii)）。

应用第 IX 讲 2.14。对\Term{set}{集合}\Term{sheaf}{层}，\Term{exact-sequence}{正合列} (4.3.1) 和第 IX 讲 2.14 表明，可取 $F=\pi_*C$，其中 $\pi:X'\to X$ \Term{finite}{有限}，$X'$ \Term{normal}{正规}因而\Term{regular}{正则}，$C$ 为 $X'$ 上的\Term{constant-sheaf}{常值层}。对挠阿贝尔\Term{sheaf}{层}，取由 $\pi_*C$ 型\Term{sheaf}{层}的\Term{finite-product}{有限积}构成的 $F$ 的一个\Term{resolution}{解消}，应用\Term{resolution}{解消}的\Term{spectral-sequence}{谱序列}，仍归结到 $F=\pi_*C$ 的情形。对\Term{ind-finite-group}{ind-有限}群\Term{sheaf}{层}，回顾，只须证明 $(R^1g_{\mathrm{cl}*})\varepsilon^*$ 对\Term{sheaf}{层} $F$ 的\Term{effaceable}{可消去}性，因而只须对可将 $F$ \Term{embedding}{嵌入}其中的\Term{sheaf}{层} $G$ 证明它，故仍归结到 $F=\pi_*C$。由第 VIII 讲 5.5 和 5.8，显然现在可将 $X$ 换成 $X'$，遂得引理。
\end{proof}
为处理最后这一情形，考虑到\Term{cohomological-dimension}{上同调维数}（第 XI 讲 7.7 及 [2] 4.14.1），显然只须证明：对\Term{complete}{完备}\Term{regular}{正则}曲线有

$(q=0)$：$X$ \Term{connected}{连通}且\Term{nonempty}{非空}，当且仅当 $X_{\mathrm{cl}}$ \Term{connected}{连通}且\Term{nonempty}{非空}。

$(q=1)$：\Term{base-change}{基变换}\Term{morphism}{态射}诱导 $X$ 的\Term{etale-covering}{平展覆盖}\Term{category}{范畴}与 $X_{\mathrm{cl}}$ 的\Term{finite}{有限}解析\Term{etale-covering}{平展覆盖}\Term{category}{范畴}的\Term{equivalence}{等价}。

$(q=2)$：有\Term{isomorphism}{同构}
\[
\varepsilon:H^2(X_{\text{ét}},\mu_n)\xrightarrow{\sim}H^2(X_{\mathrm{cl}},\mu_n).
\]
前两个断言立即由 GAGA 而得。对最后一个断言，注意 $H^2(X_{\text{ét}},\mathbb G_m)=0$（第 IX 讲 4.5）；同样，以 $\mathcal O^*$ 记
\SourcePage{160}
$X_{\mathrm{cl}}$ 上\Term{invertible}{可逆}全纯函数的\Term{sheaf}{层}，则 $H^2(X_{\mathrm{cl}},\mathcal O^*)=0$（由\Term{exact-sequence}{正合列}
\[
0\longrightarrow\mathbb Z\longrightarrow\mathcal O\xrightarrow{\exp}\mathcal O^*\longrightarrow0
\]
及 GAGA 所给出的 $H^2(X_{\mathrm{cl}},\mathcal O)=H^2(X_{\text{ét}},\mathcal O_X)=0$ 而得）。因此 \Term{kummer-exact-sequence}{Kummer 正合列}（第 IX 讲 3.2）（相应地，$X_{\mathrm{cl}}$ 上的\Term{exact-sequence}{正合列}
\[
0\longrightarrow\mu_n\longrightarrow\mathcal O^*\xrightarrow{n}\mathcal O^*\longrightarrow0
\]
）给出
\[
H^1(X_{\text{ét}},\mathbb G_m)/n\xrightarrow{\sim}H^2(X_{\text{ét}},\mu_n)
\]
（相应地，$H^1(X_{\mathrm{cl}},\mathcal O^*)/n\xrightarrow{\sim}H^2(X_{\mathrm{cl}},\mu_n)$）。由 GAGA，\Term{canonical}{典范}\Term{morphism}{态射}
\[
H^1(X_{\text{ét}},\mathbb G_m)=\operatorname{Pic}X\longrightarrow\operatorname{Pic}X_{\mathrm{cl}}=H^1(X_{\mathrm{cl}},\mathcal O^*)
\]
为\Term{bijection}{双射}，遂得结论，从而证明了 $f$ \Term{proper}{固有}时的 4.1。

\textbf{$F$ \Term{constructible}{可构造}的情形。}对 $\dim X$ 归纳。设结果在\Term{dimension}{维数} $<n$ 时已知，证明 $\dim X\leq n$ 时的结果。
\begin{SGAresult}{引理}{4.5}
可设 $f:X\to S$ 为\Term{dense}{稠密}\Term{open-immersion}{开浸入}且 $F$ 常值。
\end{SGAresult}
\begin{proof}
如 1.3 的证明，可设 $X,S$ \Term{affine}{仿射}，从而 $f$ \Term{quasi-projective}{拟射影}。应用 4.4 的论证，归结到 $F=\pi_*C$、$\pi:X'\to X$ \Term{finite}{有限}且 $C$ 常值的情形，于是将 $X$ 换成 $X'$，归结到 $F$ 常值的情形。可把 $X$ 换成 $X_{\mathrm{red}}$。取\Term{commutative-diagram}{交换图}
\[
\begin{tikzcd}
X\arrow[r,"h"]\arrow[dr,"f"']&\overline X\arrow[d,"\overline f"]\\&S
\end{tikzcd}
\]
其中 $\overline X$ \Term{reduced}{约化}，$\overline f$ \Term{projective}{射影}，$h$ 为\Term{dense}{稠密}\Term{open-immersion}{开浸入}。只须对 $h,\overline f$ 两个\Term{morphism}{态射}核实定理（参见 4.3 的证明）；$\overline f$ \Term{proper}{固有}，故只须对 $h$ 核实，亦即对\Term{dense}{稠密}\Term{open-immersion}{开浸入}核实。
\end{proof}
\begin{SGAresult}{引理}{4.6}
还可设 $S$ \Term{regular}{正则}。
\end{SGAresult}
\SourcePage{161}
\begin{proof}
设 $h:S'\to S$ 为“$S$ 的\Term{resolution-of-singularities}{奇点消解}”，亦即一个\Term{surjectivity}{满射}、\Term{proper}{固有}且\Term{birational}{双有理}的\Term{morphism}{态射}，使 $S'$ \Term{regular}{正则}；设
\[
\begin{tikzcd}
X'\arrow[r,"f'"]\arrow[d,"h'"']&S'\arrow[d,"h"]\\
X\arrow[r,"f"']&S
\end{tikzcd}
\]
为由 $h$ 得到的\Term{cartesian-square}{笛卡尔图}。

设 $F$ 为 $X$ 上的常值\Term{set}{集合}\Term{sheaf}{层}。有 $F\hookrightarrow G=h'_*h'^*F$，由此知在 $q=0$ 时只须对 $G$ 核实 4.1（参见 (4.3.1)），亦即对具有 $h'_*F'$ \Term{form}{形式}、其中 $F'$ 在 $X'$ 上常值的\Term{sheaf}{层}核实它。有\Term{commutative-diagram}{交换图}（省略“ét”\Term{signature}{符号}）
\[
\begin{tikzcd}[column sep=small]
\varepsilon^*f_*(h'_*F')=\varepsilon^*h_*f'_*F'\arrow[r,"a"]\arrow[d,"d"']&h_{\mathrm{cl}*}\varepsilon^*f_*F'\arrow[d,"b"]\\
f_{\mathrm{cl}*}\varepsilon^*(h'_*F')\arrow[r,"c"']&f_{\mathrm{cl}*}h'_{\mathrm{cl}*}\varepsilon^*F'=h_{\mathrm{cl}*}f'_{\mathrm{cl}*}\varepsilon^*F'
\end{tikzcd}
\]
由\Term{proper}{固有}\Term{morphism}{态射}的结果，$a,c$ 为\Term{bijection}{双射}。为证明 $d$ 为\Term{bijection}{双射}，只须证明 $b$ 为\Term{bijection}{双射}，亦即证明
\[
\varepsilon^*f'_*F'\xrightarrow{\sim}f'_{\mathrm{cl}*}\varepsilon^*F',
\]
故归结到 $S=S'$、从而 $S$ \Term{regular}{正则}的情形。（注记：这里未明确提及，却使用了\Term{proper}{固有}\Term{morphism}{态射}的\Term{finiteness-theorem}{有限性定理}，第 XIV 讲 1.1。）

设 $F$ 为 $X$ 上的常值挠阿贝尔\Term{sheaf}{层}。只须处理 $F=(\mathbb Z/n)_X$。\Term{morphism}{态射} $h'$ 在 $X$ 的一个\Term{dense-open-set}{稠密开集} $j:U\to X$ 上\Term{induction-isomorphism}{诱导同构}；以 $j':U'\to X'$ 记\Term{open-subset}{开集} $h'^{-1}(U)\xrightarrow{\sim}U$，令 $Y=X-U$。因 $U$ \Term{dense}{稠密}，有 $\dim Y<n$，故定理对 $(f,(\mathbb Z/n)_Y)$ 成立。由\Term{exact-sequence}{正合列}
\[
0\longrightarrow j_!(\mathbb Z/n)_U\longrightarrow(\mathbb Z/n)_X\longrightarrow(\mathbb Z/n)_Y\longrightarrow0
\]
及\Term{five-lemma}{五引理}，归结到对\Term{sheaf}{层} $j_!(\mathbb Z/n)_U$ 和\Term{morphism}{态射} $f$ 证明定理；由 4.2，也就是对 $(hf',j'_!(\mathbb Z/n)_{U'})$ 证明定理。又有\Term{exact-sequence}{正合列}
\[
0\longrightarrow j'_!(\mathbb Z/n)_{U'}\longrightarrow(\mathbb Z/n)_{X'}\longrightarrow(\mathbb Z/n)_{Y'}\longrightarrow0,
\]
且由归纳假设，定理对 $(hf',(\mathbb Z/n)_{Y'})$ 成立。因此只须对 $(hf',(\mathbb Z/n)_{X'})$ 证明定理。但由\Term{proper}{固有}情形，定理对 $(h,\cdot)$ 成立，由此经\Term{morphism}{态射}对 $h,f'$ 的 Leray \Term{spectral-sequence}{谱序列}，
\SourcePage{162}
只须对 $(f',\mathbb Z/n)$ 证明它，遂得此种情形的引理。

最后，设 $F$ 为常值\Term{finite-group}{有限群}\Term{sheaf}{层}。回顾，只须证明\Term{functor}{函子} $(R^1f_{\mathrm{cl}*})\varepsilon^*(\cdot)$ 对 $F$ \Term{effaceable}{可消去}，故只须对一个可将 $F$ \Term{embedding}{嵌入}其中的 $G$ 证明；因而可取 $G=h'_*h'^*F$，亦即具有 $h'_*F'$ \Term{form}{形式}、其中 $F'$ 在 $X'$ 上常值的\Term{sheaf}{层}。有
\[
\varepsilon^*(h'_*F')\xrightarrow{\sim}h'_{\mathrm{cl}*}(\varepsilon^*F')
\]
及
\[
0\longrightarrow(R^1f_{\mathrm{cl}*})h'_{\mathrm{cl}*}\varepsilon^*(\cdot)\longrightarrow(R^1f_{\mathrm{cl}}h'_{\mathrm{cl}*})\varepsilon^*(\cdot)=(R^1h_{\mathrm{cl}}f'_{\mathrm{cl}*})\varepsilon^*(\cdot),
\]
故只须证明后一\Term{functor}{函子}对\Term{constant-sheaf}{常值层}的\Term{effaceable}{可消去}性。但有\Term{exact-sequence}{正合列}
\[
0\longrightarrow(R^1h_{\mathrm{cl}*})f'_{\mathrm{cl}*}\varepsilon^*(\cdot)\longrightarrow(R^1h_{\mathrm{cl}}f'_{\mathrm{cl}*})\varepsilon^*(\cdot)\longrightarrow h_{\mathrm{cl}*}(R^1f'_{\mathrm{cl}*})\varepsilon^*(\cdot),
\]
因 $h$ \Term{proper}{固有}，第一项\Term{effaceable}{可消去}，例如可被 \Term{godement-resolution}{Godement 解消}消去（第 XII 讲 3.3）。所以只须证明右项被 \Term{godement-resolution}{Godement 解消}消去，即 $(R^1f'_{\mathrm{cl}*})\varepsilon^*(\cdot)$ 被该\Term{resolution}{解消}消去，亦即定理对 $(f',F')$、$F'$ 常值时成立，遂得引理。
\end{proof}
\begin{proof}[证明的结束]
现在假设 $F$ 为\Term{constant-sheaf}{常值层}，$f:X\to S$ 为\Term{dense}{稠密}\Term{open-immersion}{开浸入}，$S$ \Term{regular}{正则}。\Term{set}{集合}情形立即由 3.2 和经典\Term{topology}{拓扑}下类似的平凡引理而得：由此可知 $\varepsilon^*f_*F$ 和 $f_{\mathrm{cl}*}\varepsilon^*F$ 是取相同值的\Term{constant-sheaf}{常值层}，故经 $\varphi$ \Term{isomorphism}{同构}。

对 $q>0$，先进一步归结到 $X=S-Y$、$Y$ \Term{regular}{正则}的情形，即 $(Y,S)$ 为 $\operatorname{Spec}\mathbb C$-\Term{smooth-pair}{光滑对}的情形。事实上，写 $X=S-C_1$，其中 $C_1$ 带有诱导的\Term{reduced-structure}{约化结构}，并递归定义
\[
\begin{aligned}
Y_\nu&=C_\nu\text{ 的正则点构成的稠密开集},\\
C_{\nu+1}&=C_\nu-Y_\nu\text{，赋予诱导的约化结构}.
\end{aligned}
\]
由于 $Y_\nu$ 在 $C_\nu$ 中\Term{dense}{稠密}，有 $\dim C_{\nu+1}<\dim C_\nu$，故有序列
\[
X=X_1\subset X_2=S-C_2\subset X_3=S-C_3\subset\cdots\subset X_r=S,
\]
其中包含 $i_\nu:X_\nu\to X_{\nu+1}$ 为\Term{open-immersion}{开浸入}，且
\[
X_\nu=X_{\nu+1}-Y_\nu,\qquad Y_\nu\text{ 在 }X_{\nu+1}\text{ 中闭且正则}.
\]
$X$ 上\Term{constant-sheaf}{常值层}的\Term{direct-image}{直接像} $i_{\nu*}F$ 常值（3.2），而 $R^qi_{\nu*}F$ 对所考虑的 $q$ 值\Term{constructible}{可构造}，且集中在 $Y$ 的\Term{dimension}{维数} $\leq n-1$ 的簇上（3.4、3.6、3.7）。故由归纳假设和 Leray \Term{spectral-sequence}{谱序列}（相应地，第 XII 讲 3.2 的\Term{exact-sequence}{正合列}），归结到 $i_\nu$ 的情形，得到所声称的归约。

但对\Term{smooth-pair}{光滑对} $(Y,S)$，已在\Term{etale-topology}{平展拓扑}下明确算出 $(R^qf_*)F$ 的值（3.4、3.6、3.7）；一个类似而容易的计算，在经典\Term{topology}{拓扑}下给出类似结果，留给读者。
\SourcePage{163}
直接比较即可得到 $f$ 的定理。
\end{proof}
\SGAsection{5}{特征零代数预概形的有限性定理}
\begin{SGAresult}{定理}{5.1}
设 $k$ 为\Term{characteristic-zero}{特征零}的域，$f:X\to S$ 为在 $k$ 上\Term{locally-finite-type}{局部有限型}的\Term{scheme}{概形}之间的\Term{finite-type}{有限型}\Term{morphism}{态射}。设 $F$ 为 $X$ 上\Term{constructible}{可构造}的\Term{set}{集合}\Term{sheaf}{层}（相应地，\Term{finite-group}{有限群}\Term{sheaf}{层}；相应地，挠阿贝尔\Term{sheaf}{层}）。则 $R^qf_*F$ 在 $q=0$ 时（相应地，$q=0,1$ 时；相应地，对一切 $q$）也\Term{constructible}{可构造}。
\end{SGAresult}
假设 $k$ 的\Term{characteristic}{特征}为 $0$，是因为将使用\Term{resolution-of-singularities}{奇点消解}定理 [1]。若承认\Term{resolution-of-singularities}{奇点消解}，只要 $F$ 的阶与 $p$ 互素，该证明也适用于\Term{characteristic}{特征} $p>0$。因此应用 Abhyankar 的结果 [3]，对 $\dim X\leq2$ 可推出一个非零\Term{characteristic}{特征}的结果。除了一再用到的\Term{proper}{固有}\Term{morphism}{态射}的\Term{base-change}{基变换}定理外，证明的另一个组成部分是 3.9 的“绝对\Term{cohomology}{上同调}\Term{purity}{纯性}定理”。

回顾，对\Term{proper}{固有}\Term{morphism}{态射}及任意 $S$，定理已经证明（第 XIV 讲 1.1）。此外，在第 XIX 讲中（仍为\Term{characteristic}{特征} $0$），将对\Term{excellent}{优良}\Term{scheme}{概形}证明 3.10 所指出\Term{form}{形式}的绝对\Term{purity}{纯性}定理，从而在弱得多的假设——$S$ 为\Term{excellent}{优良}\Term{scheme}{概形}、$f$ \Term{finite-type}{有限型}——下得到\Term{finiteness-theorem}{有限性定理}。

在\Term{characteristic}{特征}非零且\Term{scheme}{概形}\Term{dimension}{维数} $\geq3$ 的情形，目前只有第 XI 讲 3.3 的如下简单推论：
\begin{SGAresult}{定理}{5.2}
设 $X$ 为 $\operatorname{Spec}k$ 上的\Term{smooth}{光滑}\Term{algebra}{代数}\Term{scheme}{概形}，$k$ 为\Term{separably-closed}{可分闭}域，$F$ 为\Term{locally-constant}{局部常值}且阶与 $k$ 的\Term{characteristic}{特征}互素的\Term{finite}{有限}\Term{abelian-group}{阿贝尔群}\Term{sheaf}{层}（相应地，\Term{finite-group}{有限群}\Term{sheaf}{层}）。则对一切 $q$（相应地，$q=0,1$），$H^q(X,F)$ 为\Term{finite-group}{有限群}（相应地，\Term{finite}{有限}\Term{pointed-set}{带基点集合}）。
\end{SGAresult}
\begin{proof}[5.2 的证明]
可设 $k$ \Term{algebraically-closed}{代数闭}（第 VIII 讲 1.1）。对 $n=\dim X$ 归纳：由第 XI 讲 3.3，存在以“\Term{good-neighborhood}{好邻域}”为\Term{open-subset}{开集}的 $X$ 的一个\Term{hypercovering}{超覆盖}。利用\Term{hypercovering}{超覆盖}的\Term{spectral-sequence}{谱序列}，立即归结到 $X$ 为\Term{good-neighborhood}{好邻域}的情形，因而存在一个\Term{elementary-fibration}{初等纤维化}（第 XI 讲 3.1）
\SourcePage{164}
\[
\begin{tikzcd}
X\arrow[r,"j"]\arrow[dr,"f"']&\overline X\arrow[d,"\overline f"]\\&S.
\end{tikzcd}
\]
由 (3.2、3.6、3.7)，\Term{sheaf}{层} $j_*F$ 在 $X$ 上\Term{locally-constant}{局部常值}，$(R^1j_*)F$ 在 $Y=\overline X-X$ 上\Term{locally-constant}{局部常值}，且当 $q>0$ 时 $(R^qj_*)F=0$。因此由 2.1，对所考虑的 $q$ 值，$(R^p\overline f_*)(R^qj_*)F$ 是 $S$ 上\Term{locally-constant}{局部常值}的\Term{sheaf}{层}；于是由归纳假设和 Leray \Term{spectral-sequence}{谱序列}（相应地，第 XII 讲 3.2 的\Term{exact-sequence}{正合列}）得出结论。
\end{proof}
\begin{proof}[5.1 的证明]
该证明与 $F$ \Term{constructible}{可构造}时 4.1 的证明极为接近。这里只给出大纲：对 $\dim X=n$ 归纳。先用 1.3 证明的方法，归结到 $X,S$ \Term{affine}{仿射}、从而 $f$ \Term{quasi-projective}{拟射影}的情形。再用第 IX 讲 2.15 及第 VIII 讲 5.5 和 5.8，归结到 $F$ 常值的情形。有\Term{commutative-diagram}{交换图}
\[
\begin{tikzcd}
X\arrow[r,"h"]\arrow[dr,"f"']&\overline X\arrow[d,"\overline f"]\\&S,
\end{tikzcd}
\]
其中 $h$ 为\Term{dense}{稠密}\Term{open-immersion}{开浸入}，$\overline f$ \Term{proper}{固有}，故定理对 $\overline f$ 成立（第 XIV 讲 1.1）。因此只须对 $(h,F)$ 证明定理，亦即对\Term{dense}{稠密}\Term{open-immersion}{开浸入}和\Term{constant-sheaf}{常值层}证明它。设 $h:S'\to S$ 为 $S$ 的\Term{resolution-of-singularities}{奇点消解}，且
\[
\begin{tikzcd}
X\arrow[d,"f"']&X'\arrow[l,"h'"']\arrow[d,"f'"]\\
S&S'\arrow[l,"h"]
\end{tikzcd}
\]
为由 $h$ 得到的\Term{cartesian-square}{笛卡尔图}。归结到对 $f'$ 证明定理，即进一步假设 $S=S'$ \Term{regular}{正则}的情形。用 4.1 证明末尾的归纳方法，归结到 $Y=S-X$ \Term{nonsingular}{非奇异}、亦即 $(Y,S)$ 为\Term{smooth-pair}{光滑对}的情形，最后应用 (3.2、3.4、3.6、3.7) 即告完成。
\end{proof}
\textbf{注记 5.3。}刚才概述的 5.1 的证明还适用于以下情形：给定一个被整数 $n>0$ 零化的左\Term{noetherian}{诺特}环 $A$（若不假设\Term{characteristic}{特征}为零，就必须假设 $n$ 与\Term{characteristic}{特征}互素，
\SourcePage{165}
参见 5.1 后的注记），并考虑\Term{constructible}{可构造}左 $A$-模\Term{sheaf}{层}。

\begin{center}\textbf{参考文献}\end{center}
\begin{enumerate}
\item[{[1]}] H. Hironaka，\textit{Resolution of singularities of an algebraic variety over a field of characteristic zero}（\Term{characteristic-zero}{特征零}域上\Term{algebraic-variety}{代数簇}的\Term{resolution-singularities}{奇点消解}），Annals of Math.，第 79 卷（1964），第 109 页。
\item[{[2]}] P. Godement，\textit{Théorie des Faisceaux}（层论），巴黎，1958。
\item[{[3]}] S. Abhyankar，\textit{Local uniformization on algebraic surfaces over ground fields of characteristic $p\ne0$}（\Term{characteristic}{特征} $p\ne0$ 的\Term{base-field}{基域}上\Term{algebra}{代数}曲面的\Term{local-scheme}{局部}一致化），Ann. of Math.，第 63 卷（1956），第 491 页。
\item[{[4]}] J. P. Serre，\textit{Géométrie Algébrique et Géométrie Analytique}（\Term{algebra}{代数}几何与解析几何），Annales Inst. Fourier，1956，第 1--42 页（引用时简称 GAGA）。
\end{enumerate}
